% Zero-Shot Self-Orchestration with Ledger-Based Control for Improved LLM Coding Performance
% Converted from "PAPER OUTLINE copy.md".
%
% Build:  pdflatex paper.tex && pdflatex paper.tex     (twice, for the ToC-free refs/floats)
%   or:   latexmk -pdf paper.tex
%   or:   tectonic paper.tex
%
% Every figure is read live from ../escalation/plots/; nothing is copied into this
% directory. Regenerate with, in escalation/run_bench_script/:
%   plot_4new_5pass_reason_on.py    -> Figure 1, panel (a)
%   plot_cost_5_pass.py             -> Figure 2, a single panel over a stats table
%   plot_cost_vs_score.py           -> Figure 3, cost against accuracy
%   plot_16k_reason_off_5_pass.py, plot_128k_reason_off_1_pass.py,
%   plot_128k_reason_on_1_pass.py   -> Figures 4-6, panel (a)
%   plot_bars.py                    -> Figures 1-2, 4-6, panel (b)
%   make_figures_tex.py             -> fig-*.tex, the \input files below: both panels
%                                      plus a caption whose notes are generated from
%                                      the same stats the charts plot. Edit those notes
%                                      in the plot script, never in fig-*.tex.
%   plot_agent_loop_flowchart.py    -> Figure 7
% The old ./media/ directory (pandoc extracts from the Markdown draft) is gone; it is
% in history at ff4313c if any of those bitmaps are ever needed again.
%
% Section cross-references are written literally (\S2.3, \S3.1) rather than with
% \label/\ref: the numbering here is fixed and already matches the prose, and literal
% refs cannot degrade to "??" if a label is ever dropped.

\documentclass[11pt]{article}

\usepackage[utf8]{inputenc}
\usepackage[T1]{fontenc}
\usepackage{lmodern}          % scalable outline font; microtype expansion needs one
\usepackage[letterpaper,margin=1in]{geometry}
\usepackage{amsmath}
\usepackage{amssymb}
\usepackage{booktabs}
\usepackage{graphicx}
\usepackage{array}
\usepackage{textcomp}
\usepackage{caption}
\usepackage{subcaption}
% \FloatBarrier only. The figures are [p] float pages, which LaTeX is free to carry
% past a \subsection heading -- a chart discussed in one subsection then prints inside
% the next. A barrier at the end of a subsection that owns a figure pins it there.
\usepackage{placeins}
\usepackage{microtype}
\usepackage[hidelinks]{hyperref}

% Block paragraphs: no first-line indent, a visible gap between paragraphs
% instead. parskip is loaded rather than setting the lengths by hand because it
% also keeps list and float spacing from inheriting the new \parskip.
\usepackage{parskip}

% Charts are read straight out of the generating directory, so a re-run of
% escalation/run_bench_script/plot_*.py updates the paper with no copy step.
% Panel (a) of each figure is the dot/line chart, panel (b) the bar chart of the
% same numbers -- both are emitted by those scripts from one shared compute().
% BUNDLE EDIT (the only change from the repository copy): the figures ship in
% ./plots beside this file rather than in the repo's ../escalation/plots.
\newcommand{\plotdir}{plots}
% Panels are drawn at \linewidth, and that is the bound that has to bind: each chart is
% authored on a canvas exactly twice its printed size (6.5 x 3.1in), so its type lands
% on the page at the sizes named in the plot scripts -- a title that reads as \large, 10.95pt labels,
% 9pt fine print, matching \large, \normalsize and \footnotesize here. A height cap that
% bit would scale the image down and shrink all of that with it, so the cap below sits
% just above the panels' own 3.1in and only catches a chart regenerated much taller.
% 2 x 3.1in of image leaves ~2.8in for the three captions, which the notes need. LaTeX
% does not flag a float page that overruns: check the last caption line on the PDF.
\newcommand{\panelplot}[1]{%
  \includegraphics[width=\linewidth,height=0.36\textheight,keepaspectratio]{#1}}

\captionsetup[table]{skip=6pt}
\captionsetup{font=small,labelfont=bf}

% A float is allowed its own page only once it fills three-quarters of one. The default
% half lets a single-panel figure -- chart plus a long caption, about two-thirds of a page
% -- claim a page and leave the rest white; at 0.75 it goes to the top of a text page and
% the text follows it down.
\renewcommand{\floatpagefraction}{0.75}

\title{\bfseries Zero-Shot Self-Orchestration with Ledger-Based Control for Improved LLM Coding Performance}

% \maketitle switches the title footnotes to article's symbol sequence (*, dagger,
% double dagger). The three marks below are an equal-contribution note, an
% affiliation and the corresponding author, so they read better as a, 1, dagger.
% There are no other footnotes in the document, so overriding the sequence itself
% is enough.
\makeatletter
\renewcommand*{\@fnsymbol}[1]{\ifcase#1\or a\or 1\or \ensuremath{\dagger}\else
  \@ctrerr\fi}
\makeatother

% The whole list on one line is wider than the text block, and article's \author is a
% tabular cell that cannot break: it would run past the right margin and read as
% off-centre. Hence the explicit \\ -- each row is centred on its own, and the break
% falls between authors rather than inside a name.
\author{%
  Victor Gao\thanks{Persis Capital Inc.}\hspace{0.35em}\thanks{These authors contributed equally to this work.}\quad
  Vida Khosrowshahi\footnotemark[1]\hspace{0.35em}\footnotemark[2]\quad
  Ali Khosrowshahi\footnotemark[1]\hspace{0.35em}\footnotemark[2]\\[0.4ex]
  Xihao Sun\footnotemark[1]\quad
  Juhyun Lee\quad
  Simon (Sang Won) Lee\footnotemark[1]\hspace{0.35em}%
  \thanks{Corresponding author: Simon (Sang Won) Lee, Ph.D.
     \href{mailto:slee@persisholdings.com}{slee@persisholdings.com}.}
}
\date{}

\begin{document}
\maketitle

\begin{abstract}
% abstract is a quotation environment, whose own 1.5em indent parskip does not reach.
\noindent
Frontier coding performance is often pursued by using larger proprietary models at
substantial inference cost. We introduce ledger-based zero-shot self-orchestration, a
training-free manager–worker method in which fresh instances of the same model decompose
problems, verify candidate solutions, and coordinate through a shared filesystem
containing a plan, task list, notes, and current solution. We compare the method with
single-call baselines across nine open- and closed-weight models on the 100 latest
\textit{hard} LiveCodeBench problems. On pinned backends with reasoning enabled and a
128k output cap, five independent passes show gains of 23.4 percentage points for
Qwen3.8-27B, 10.6 for GPT-5.6-Luna, and 8.0 for GPT-5.6-Terra. The method provides two
routes to Fable 5-level coding performance. Managed GPT-5.6-Terra achieves 85.0\% mean
pass@1 versus Fable 5's 87.4\% single-call result at \$11.71 versus \$61.11 per
100-problem pass---19\% of the cost; the 2.4-point accuracy gap is unresolved by a paired
permutation test over problems (Holm-corrected $p = 0.59$), with a one-sided 95\%
confidence bound admitting a deficit of up to 5.8 points. Separately, locally served,
open-weight Qwen3.8-27B rises from 63.0\% to 86.4\%, one point below Fable 5; this gap is
likewise unresolved ($p = 0.73$), with the corresponding bound admitting a deficit of up
to 4.8 points. Transcript analysis identifies problem decomposition and persistent
context management as recurring mechanisms, although effects remain model-dependent: for
Qwen3.6-35B with reasoning disabled, the effect is statistically null over five passes at
a 16k cap ($-1.2$ points; Holm-corrected $p = 0.65$) and $-9$ points in a single 128k
pass. These results show that inference-time organization can approach frontier
single-pass coding accuracy either at substantially lower API cost or with a
self-hostable open model.
\end{abstract}

\section{Introduction}

Large Language Models (LLMs) are increasingly deployed not as single-shot predictors but as
\textit{agents}: models wrapped in a scaffold that lets them plan, call tools, keep notes,
and revise their own work over multiple turns. A popular next step is to use a system of
\textit{several} such agents together, on the intuition that a team that decomposes a
problem, critiques one another, and pools partial results should outperform a single model
answering in a single pass. That intuition has produced a large literature
\cite{guo2024}, but the evidence is mixed \cite{wangq2024} and much of it is
confounded \cite{tran2026}: multi-agent pipelines usually change
several factors at once, including token budgets, tool calls, prompts, and retrieval, so
an aggregate gain over a single call rarely identifies which factor helped.

Our comparison holds the underlying model and problem set fixed while 
comparing two conditions: the same model answering in a single pass and the same model 
operating within a manager-plus-workers scaffold over a shared filesystem workspace. 
The scaffold requires additional test-time computation but no training or per-benchmark
tuning. Its coordination is dynamic rather than learned: the manager revises the task list
after each round, chooses the next step, and decides when to stop. Every role uses the same model in
a fresh context.

\subsection{Related work}

We group prior multi-agent LLM systems into four families according to how they
coordinate, and then compare our design with each family. These systems build on single-agent
scaffolds that add structure around one model: \textit{ReAct} \cite{yao2022}
interleaves reasoning traces with tool actions, \textit{Reflexion}
\cite{shinn2023} has an agent verbally critique its own attempt and
retry from that reflection, and \textit{Self-Refine} \cite{madaan2023} iterates
generate$\to$feedback$\to$revise within a single model. Our
worker role is an instance of that idea; what we vary is the coordination layer above it.

\paragraph{(a) Learned / trained orchestrators.}
Here the coordination policy itself is \textit{trained}: a model learns to decompose a
task, assemble a team, and route or role-assign sub-tasks. It is the line closest to what
we do. Sakana AI's Fugu \cite{sakana2026}
packages a whole multi-agent orchestration system behind a single model
API - a learned coordinator that, per query, assembles and coordinates a pool of expert
LLM workers, verifies, and synthesises. Its authors position it against two alternatives:
multi-agent workflows that users must design, tune and operate, and coordination built on
fixed communication patterns or single-step routing - Fugu instead trains the
orchestrator itself to decide, adaptively and per query, how to use its agent pool. The
more general of its two ICLR 2026 foundations is the Conductor
\cite{nielsen2025}, a 7B model RL-trained (GRPO) to \textit{design worker communication
topologies} and per-worker prompts - reaching state of the art on GPQA-Diamond and
LiveCodeBench \textit{by orchestrating other models rather than solving problems itself}.
What is learned here is the workflow structure itself. (Its companion, Trinity,
learns only role \textit{assignment} over a fixed structure, and we discuss it under (b).)
Earlier learned-combination work includes query routers such as \textit{RouteLLM}
\cite{ong2024}, which trains a router on preference data to pick a model per query.
Our system is the training-free counterpart to this family:
the manager is a fixed prompt with no learning, which lets us ask how much of the
orchestration benefit is available \textit{without} training, and \textit{when} it
materialises.

\paragraph{(b) Fixed task lists and hand-designed workflows.}
Roles and their hand-off order are specified in advance and every problem flows through
the same pipeline. \textit{MetaGPT} \cite{hong2023} encodes a
software company's SOPs across PM/architect/engineer/QA roles; \textit{ChatDev}
\cite{qian2023} runs a chat-driven design$\to$code$\to$test waterfall;
\textit{AutoGen} \cite{wu2023} provides configurable conversable
agents; \textit{CAMEL} \cite{li2023} fixes a role-playing pair. For
code specifically, \textit{AgentCoder} \cite{huang2023} pairs a
coder with test-designer and test-executor roles in a fixed loop, and \textit{MapCoder}
\cite{islam2024} runs a fixed four-agent pipeline - example
recall, planning, code generation, debugging. Trinity \cite{xu2025} sits at the
boundary of this family and the learned one: it fixes a
three-role workflow - Thinker, Worker, Verifier - but uses a small
evolution-strategy-trained coordinator to decide, each turn, which LLM fills each of the
three fixed roles. The \textit{structure} is hand-designed; only the per-turn model
\textit{assignment} is learned. This family is closest to ours in spirit, but its pipeline
is static, whereas our manager re-curates the task list every step and decides for itself
when to stop.

\paragraph{(c) Shared-blackboard and decentralised contributions.}
Here agents coordinate through a shared external store or aggregation rule rather than a
hard-coded pipeline. Most relevant to us, ARIADNE \cite{wei2026} targets exactly
our setting, program generation, with a
blackboard-driven MCTS: candidate solutions are explored under a reward signal,
with a shared blackboard as the working memory that the search reads and writes.
LbMAS \cite{han2025} likewise mediates all agent
communication through a global blackboard, choosing agents dynamically from board
state rather than a fixed template, which also cuts the overall prompt length across
agents. In the fully
decentralised direction, \textit{multi-agent debate} \cite{du2023} has agents critique
and revise toward consensus (a ``society of minds''), and
\textit{sampling-and-voting} \cite{li2024} shows raw agent count alone can lift accuracy;
see Guo et al.\ \cite{guo2024} for a survey. \textit{Mixture-of-Agents} \cite{wangj2024a}
aggregates in layers rather than in one vote: several proposer models answer, and an
aggregator model synthesises their answers into the next layer's input. As a zero-shot orchestrator,
nothing is trained - the aggregator is an off-the-shelf model given a fixed synthesis prompt, and
which model sits in which layer is chosen by hand from win rate and output diversity.
Our workspace is a shared store that keeps each call's
context short, similar to the blackboard used in ARIADNE/LbMAS, but our coordination is a
single manager's plain task-list loop, with no MCTS, no reward model, and no
learned search policy.

\paragraph{\textit{(d) Budget-controlled comparisons.}}
A separate line of work asks not how coordination is decided but whether reported
multi-agent gains survive equalising test-time compute. Wang et al.\ \cite{wangj2024b}
price each strategy in tokens and compare at a matched budget: several of the more
elaborate ones keep little of their advantage over plain self-consistency once the spend
is held equal, which is their argument for reporting budget alongside accuracy as a matter
of course. Tran and Kiela \cite{tran2026}
match the intermediate reasoning-token budget between a single agent and
several multi-agent architectures across Qwen3, DeepSeek-R1-Distill-Llama and Gemini 2.5,
and find that single-agent systems match or outperform multi-agent ones on multi-hop
reasoning once thinking tokens are held constant; they motivate this with a Data
Processing Inequality argument - routing information through additional agents cannot
add information - and predict that ``multi-agent systems become competitive when a
single agent's effective context utilization is degraded, or when more compute is
expended.'' Our experiment does not attempt an equal-token comparison on the same model: the
manager--worker loop necessarily spends more tokens than the single-agent baseline, and we
ask whether that additional spend buys better solutions - not whether orchestration is
more information-efficient per token. We also compare the cost-effectiveness of using this scaffold on a cheaper model against fewer tokens on an expensive model.

\subsection{Our experiment}

Throughout, we use \textit{agent} to mean a separately prompted model invocation with a
distinct role and a fresh context; all agents share the same underlying model. We call this 
zero-shot self-orchestration: inference-time orchestration in which the orchestrator is 
neither trained for orchestration nor provided task-specific demonstrations of how to 
decompose or coordinate the problem.

Our design, sits between the three families:

\begin{itemize}
  \item \textbf{A shared filesystem workspace.} The workspace serves as a
  persistent shared ledger, containing a plan, a task list, an accumulating
  notes file, and the current best solution. All roles read and write these
  files, so the state of the computation persists across agent invocations
  rather than residing in any one context window.

  \item \textbf{A manager that adapts the plan.} A \textit{manager} instance reads the
  problem, writes an overarching plan, and then runs a loop: it inspects progress, curates
  the task list, spawns a fresh \textit{worker} to do the single most valuable next task,
  and verifies the output against the sample cases (in the v2 scaffold, \S3), repeating
  until it judges the problem solved or a small round budget is exhausted. There is no
  fixed pipeline - the manager decides, per problem, what happens next.

  \item \textbf{No training, no per-benchmark tuning.} Every role is the \textit{same}
  model invoked in a fresh context with a short generic prompt. As a zero-shot orchestrator,
  nothing is trained or hand-tuned against the problem set. The two conditions hold the model, the benchmark and
  the solver temperature fixed; the ``manager'' condition adds the orchestration calls,
  their role-specific prompts, the shared workspace state and the test-time compute that
  comes with them.
\end{itemize}

This design measures the effect of the manager--worker scaffold relative to the
corresponding single-call baseline while holding the model and problem set fixed.
\section{Results}

\subsection{Pinned-backend arms at 128k with thinking on - five passes}

Our headline condition is four models run five independent times each at a 128k cap with
thinking on, on the 100 latest \textit{hard} LiveCodeBench problems (LCB-100, \S3.2).
Every arm is served on a pinned backend - the GPT-5.6 pair by the official
OpenAI API, Qwen3.8-27B by our own vLLM, Claude Fable 5 by the Anthropic Messages API -
and every arm runs the v2 scaffold (\S3). Every condition here is repeated five times. The
earlier runs of \S2.4 cover more models, but only their 16k reasoning-off condition repeats;
the rest are one pass each, through a gateway whose routing proved to be the main source of
run-to-run noise.

\begin{table}[htbp]
\centering
\caption{\textbf{LCB-100 pass@1 (\%), five independent passes.} At a 128k cap with thinking
on, all arms served on a pinned backend and run on the v2 scaffold (\S3). Mean $\pm$ standard
deviation (SD) across passes;
$\Delta$ is the paired per-pass difference (manager $-$ single).}
\small
\begin{tabular}{@{}llcccl@{}}
\toprule
\textbf{Model} & \textbf{Serving} & \textbf{Single} & \textbf{Manager} & \textbf{$\Delta$} &
\textbf{Per-pass $\Delta$} \\
\midrule
Claude Fable 5 & Anthropic & $\mathbf{87.4 \pm 1.1}$ & -\textsuperscript{b} & - & -\textsuperscript{b} \\
GPT-5.6-Terra & OpenAI & $77.0 \pm 1.0$ & $\mathbf{85.0 \pm 1.0}$ & $\mathbf{+8.0 \pm 0.0}$\textsuperscript{c} & $+8, +8, +8, +8, +8$ \\
GPT-5.6-Luna  & OpenAI & $67.2 \pm 4.3$ & $\mathbf{77.8 \pm 2.0}$ & $\mathbf{+10.6 \pm 5.1}$ & $+17, +7, +13, +4, +12$ \\
Qwen3.8-27B   & local vLLM & $63.0 \pm 4.1$\textsuperscript{a} & $\mathbf{86.4 \pm 2.7}$ & $\mathbf{+23.4 \pm 6.6}$ & $+15, +20, +29, +22, +31$ \\
\bottomrule
\end{tabular}

\vspace{0.6ex}
\begin{minipage}{0.92\linewidth}
\footnotesize
\textsuperscript{a}~Qwen3.8-27B's manager arm ran at 128k, but its single arm was
generated at a 250k cap. The
column reports that arm cap-matched back to 128k so the row is like-for-like; the procedure
is described in \S3.2. As generated at 250k it scores $65.6 \pm 4.6$ for
$\Delta = +20.8 \pm 7.0$ (\S2.3).\\
\textsuperscript{b}~Fable 5 was run single-only, so it contributes a single-call figure and
no $\Delta$. Its per-pass scores are 86, 87, 87, 88, 89.\\
\textsuperscript{c}~The zero SD is a coincidence of aggregates, not a fixed set of problems.
Across the five passes the manager wins 9/10/10/11/11 problems the single call loses and
loses 1/2/2/3/3 that it wins; the two move together and happen to net to $+8$ each time.
The underlying sets churn: 29 distinct problems are manager-only in at least one pass, and
only one is manager-only in all five.
\end{minipage}
\end{table}

\textbf{All three models evaluated in both conditions benefit from the manager}, and the
five repeated passes show that the gains are not single-pass artifacts. Terra, already strong
single-shot at 77.0, moves $+8.0$.
Luna gains $+10.6 \pm 5.1$, and the manager also halves its run-to-run spread
($4.3 \to 2.0$ SD).
Luna's per-pass $\Delta$ ranges from $+4$ to $+17$, so the effect size is real but the
per-pass estimate is unstable; Terra's narrow band is the more reliable of the two.

\textbf{Qwen3.8-27B shows the largest gain of the three}, $+23.4 \pm 6.6$ ($+15$ at
worst). The manager arm reaches $86.4 \pm 2.7$ - comparable to Fable 5 and Opus-5's (\S2.4) single-call scores -
and it does so while narrowing the spread from $4.1 \to 2.7$ SD. The direction matches the
reasoning-on pattern of \S2.4: the smaller the model, the more the scaffold has to add.

Figure 1 compares the four models. Ordered by single-call score, the manager's gain
shrinks monotonically as the single call gets stronger - $+23.4$, $+10.6$, $+8.0$ - and
Fable 5's lone bar marks where that trend is heading: the scaffold buys most where the
model unaided is weakest, and the three managed arms converge into a band a few points
below the best single call in the study rather than passing it.

% Generated by paper_plot_script/make_figures_tex.py -- do not edit.
% The caption text lives in that chart's plot script, next to the numbers it
% describes; re-run the script to update this file.
\begin{chartfigure}
\centering
\halfplot{\plotdir/manager_vs_single_call_four_models_lcb-100_5_passes_128k_max_tokens_reasoning_on_bars_light.pdf}
\caption{\textbf{Manager vs.\ single call.} 128k $\times$ 5 passes, reasoning on. Bars are pass@1 on the same 100 problems, the line through each the 95\% CI across the 5 passes (t, df = 4). The top bar of each pair is the manager, dark; the single call is under it, light; hatch is the model. ``Single call'' is one call with no tools and no loop; Fable 5 ran single-only, so it has one bar, which the dashed rule carries across the chart. The bracket beside each pair is manager $-$ single call; the right-hand column tests the top bar against Fable 5. Both are paired sign-flip permutation tests, unit = problem (n = 100), Holm-corrected within each family of 3: * p $<$ .05, ** p $<$ .01, *** p $<$ .001. The three within-model $\Delta$ all clear p $<$ 1e-4; against Fable 5, p = GPT-5.6-Terra 0.59, GPT-5.6-Luna 4.6$\times$10$^{-4}$, Qwen3.8-27B 0.73. Manager-only vs single-only problem-passes (500 per model): GPT-5.6-Terra 51/11; GPT-5.6-Luna 71/18; Qwen3.8-27B 125/8; exact McNemar 3$\times$10$^{-7}$ or better. The gains are not compensating wins and losses. Every arm is at a 128k cap and re-scored on the corrected evaluator (Appendix C); Qwen3.8-27B's single arm is the 128k cap-matched replay of a 250k generation, which Appendix B describes and Section 5.2 prices.}
\label{fig:main}
\end{chartfigure}


\paragraph{Cap-matching.}
The arms are not natively cap-matched. The GPT-5.6 pair ran both arms at 128k, so did
Fable 5's single arm and Qwen3.8's manager arm, but Qwen3.8's single arm ran at 250k, and
it spent that budget: its one call per problem hit \texttt{finish\_reason=length} at
250,000 tokens on 124 of 500 problem-passes. We therefore replay that arm at 128k by the
procedure in \S3.2, which cuts 150 of the 500 generations. What the cap is worth, and how
much of the $\Delta$ turns on it, is quantified in \S2.3.

\paragraph{Per-problem agreement.}
Pass@1 differences can hide compensating gains and losses, so we also count, over every
problem $\times$ pass, where exactly one arm succeeded. Fable 5 is absent here for want of
a second arm:

\begin{table}[htbp]
\centering
\caption{\textbf{Per-problem agreement between the two arms.} Counted over every problem
$\times$ pass, 500 per model.}
\small
\begin{tabular}{@{}lccccc@{}}
\toprule
\textbf{Model} & \textbf{Manager only} & \textbf{Single only} & \textbf{Both} &
\textbf{Neither} & \textbf{Exact McNemar \textit{p}} \\
\midrule
Qwen3.8-27B\textsuperscript{a} & \textbf{125} & 8 & 307 & 60  & $\mathbf{4 \times 10^{-28}}$ \\
GPT-5.6-Luna  & \textbf{71} & 18 & 318 & 93  & $\mathbf{1 \times 10^{-8}}$ \\
GPT-5.6-Terra & \textbf{51} & 11 & 374 & 64  & $\mathbf{3 \times 10^{-7}}$ \\
\bottomrule
\end{tabular}

\vspace{0.6ex}
\begin{minipage}{0.85\linewidth}
\footnotesize
\textsuperscript{a}~Counted against Qwen3.8's cap-matched 128k single arm, matching its row
in \S2.1. Against the as-generated 250k arm the split is 114 / 10 / 318 / 58
($p = 2 \times 10^{-23}$): cap-matching moves twelve problem-passes from ``both'' to
``manager only'' and two from ``single only'' to ``neither'', with one moving back - on
arc194\_e the 128k cut lands inside the reasoning, where the fallback of \S3.2 finds a
whole program that the 250k answer had been cut off part-way through.
\end{minipage} 
\end{table}

All three clear significance comfortably: in every case the manager wins around four times
as many problem-passes as it loses, or better. Qwen3.8's split is the most lopsided in the study at
125 against 8 - nearly sixteen to one, a net of 117 problem-passes, which is the whole $+23.4$.
Luna's is the closest, 71 against 18, and Terra's is the smallest in absolute terms, 51
against 11 for a net of 40.

% Table 2 belongs to \S2.1 and is read by the paragraph above it. Without this it floats
% past the \S2.2 heading and prints under a section that never mentions it.
\FloatBarrier

\subsection{What the scaffold costs}

Every $\Delta$ in \S2.1 is bought with tokens. The manager replaces one call with a plan,
a brainstorm, a verifier and up to ten worker rounds, and every one of those is billed:
$+153\%$ on Qwen3.8-27B (\$20.44 to \$51.75 a pass), $+266\%$ on GPT-5.6-Luna
(\$0.41 to \$1.50) and $+244\%$ on GPT-5.6-Terra (\$3.41 to \$11.71), every one of them
clear per run and per problem. Roughly, the scaffold triples the bill, and that is what
the gains in \S2.1 cost. Every arm priced here is at the same 128k output cap, Qwen3.8-27B's
single call included (\S3.2).

% Both cost figures go here, BEFORE every paragraph that reads off them, and the barrier
% below is what makes that true on the page rather than only in the source: without it
% LaTeX carries them past the readings and the reader meets "Figure 3 plots..." pages
% before Figure 3. The cost is a short page here, where only the opening paragraph
% precedes the flush; the alternative was figures the text forward-refers to for two
% pages. The barrier also pins both inside \S2.2.
%
% The order is tables, chart, chart -- the two cost tables come from fig-cost's plot
% script but are \input first on purpose. The barrier flushes floats in declaration order,
% so this is what decides which sheet each one lands on: sent last they trailed a page
% behind the charts they belong with.
% Generated by paper_plot_script/make_figures_tex.py -- do not edit.
% The caption text and the table numbers live in that chart's plot script;
% re-run the script to update this file.
\begin{table}[!hb]
\centering
\small
\caption{\textbf{What one pass cost each arm.} List rate $\times$ the tokens it consumed.}
\label{tab:cost}
\begin{tabular}{@{}l@{\hspace{0.7em}}l@{\hspace{0.7em}}r@{\hspace{0.7em}}r@{\hspace{0.7em}}r@{\hspace{0.7em}}r@{}}
\toprule
\textbf{Arm} & \textbf{Rate \$/MTok in / out} & \textbf{In (MTok)} & \textbf{Out (MTok)} & \textbf{\$/pass} & \textbf{\$/solved} \\
\midrule
Qwen3.8-27B single & \$0.35 / \$2.75 & 0.0753 & 7.4247 & \$20.44 & \$0.32 \\
Qwen3.8-27B manager & \$0.35 / \$2.75 & 1.5053 & 18.6277 & \$51.75 & \$0.60 \\
GPT-5.6-Luna single & \$0.20 / \$1.20 & 0.0661 & 0.3299 & \$0.41 & \$0.006 \\
GPT-5.6-Luna manager & \$0.20 / \$1.20 & 1.1686 & 1.0522 & \$1.50 & \$0.019 \\
GPT-5.6-Terra single & \$2 / \$12 & 0.0661 & 0.2728 & \$3.41 & \$0.044 \\
GPT-5.6-Terra manager & \$2 / \$12 & 1.1098 & 0.7911 & \$11.71 & \$0.14 \\
Fable 5 single & \$10 / \$50 & 0.0899 & 1.2043 & \$61.11 & \$0.70 \\
\bottomrule
\end{tabular}
\par\vspace{0.6ex}
{\footnotesize List rates: Qwen3.8-27B \citep{openrouter2026}, GPT-5.6-Luna and GPT-5.6-Terra \citep{openai2026price}, Fable~5 \citep{anthropic2026price}. OpenAI rates are the short-context tier and Anthropic's the base rates; no long-context tier, Batch discount or prompt-caching multiplier is applied.}
\end{table}


% Generated by paper_plot_script/make_figures_tex.py -- do not edit.
% The caption text and the table numbers live in that chart's plot script;
% re-run the script to update this file.
\begin{chartfigure}
\centering
\halfplot{\plotdir/what_one_pass_costs_lcb-100_5_passes_single_call_vs_manager_against_fable_5_light.pdf}
\caption{\textbf{The scaffold's bill.} Cost of one pass over the same 100 problems; Table~\ref{tab:cost} is the arithmetic behind every bar. No cached-input discount is taken, and Qwen3.8-27B is priced at OpenRouter market rates. The top bar of each pair is the manager, dark; the single call is under it, light; all are at a 128k output cap, and hatch is the model. The right-hand column tests each top bar against Fable 5's single call (Welch per run, $n=5$ vs 5; Holm-corrected across seven comparisons): * $p<.05$, ** $p<.01$, *** $p<.001$. Unbracketed comparisons: each manager $-$ single increase is significant ($8.7\times10^{-5}$ or better per run, $<3.5\times10^{-5}$ per problem); GPT-5.6-Luna's manager undercuts GPT-5.6-Terra's single call by \$1.91 ($2.1\times10^{-7}$ per run, $<3.5\times10^{-5}$ per problem). At 1.18x the assumed Qwen rate (\$0.413/\$3.25 per MTok, inside the spread across hosted providers) the manager's cost advantage over Fable 5 disappears entirely — a larger uncertainty than any p-value here.}
\label{fig:cost}
\end{chartfigure}


% Generated by paper_plot_script/make_figures_tex.py -- do not edit.
% The caption text lives in that chart's plot script, next to the numbers it
% describes; re-run the script to update this file.
\begin{figure}[htbp]
\centering
\panelplot{\plotdir/what_accuracy_costs_lcb-100_5_passes_128k_max_tokens_reasoning_on_all_seven_pinned-backend_arms_light.pdf}
\caption{\textbf{Cost against accuracy.} One point per arm; an arrow runs from the single call to the manager of each model that has both. x is dollars for one pass over the 100 problems (log scale), y is pass@1; the bar through each point is the 95\% CI across the 5 passes (t, df = 4). Light fill = single call, dark = with manager; marker shape is the model. Qwen3.8-27B's single arm is the 128k cap-matched one on BOTH axes - score from the replay, output tokens capped at 128,000 per call to match. Priced as generated at 250k it would sit at \$29.08 rather than \$20.44. Retried and discarded attempts are counted: they were generated and would be billed. The cheapest arm is GPT-5.6-Luna, single call at \$0.41 a pass and the most accurate is Fable 5, single call at \$61.11 - a 149$\times$ spread in price for +20.2 points.}
\label{fig:cost-vs-score}
\end{figure}


\FloatBarrier

Figure 3 plots price against accuracy for all seven arms and connects each model's two
conditions. The joint view reveals three patterns that neither measure shows alone.

\textbf{The base cost varies considerably across models}.
A single-call pass generates 7.4M output tokens on Qwen3.8-27B against 0.33M
on GPT-5.6-Luna and 0.27M on GPT-5.6-Terra. The open model thinks its way to $63.0\%$ at
length; the OpenAI models reach $67.2\%$ and $77.0\%$ in a twentieth of the tokens. Since
thinking bills as output, a bill here is as much a fact about how much reasoning a provider
lets a model emit as about the scaffold wrapped around it.

\textbf{The scaffold's move is up and to the right in every case}. It buys accuracy with money: 9.7 points per extra dollar on GPT-5.6-Luna,
1.0 on GPT-5.6-Terra, 0.7 on Qwen3.8-27B.

\textbf{Costs vary widely, while the frontier is nearly flat at the top.} The cheapest arm and
the most accurate are 149$\times$ apart in price for 20.2 points, and the last 2.4 of those
points cost more than the first 17.8.

\paragraph{Three comparisons along the frontier.}
First, GPT-5.6-Terra with a manager nearly matches Fable 5 for a fifth
of the price: $85.0$ against $87.4$, at \$11.71 a pass against \$61.11. Those 2.4 points
are within error ($p = 0.59$, Figure 1), and the one-sided 95\% bound allows a
deficit of 5.8 points, while the savings are significant ($p < 10^{-4}$ both per run and per problem).
Both arms are priced off API list rates in effect when the experiments were run
\cite{openai2026price,anthropic2026price}.

Second, Qwen3.8-27B with a manager reaches $86.4$ against Fable 5's $87.4$, and
does it for less: \$51.75 a pass against \$61.11, a \$9.36 saving per 100-problem pass
($p = 0.005$). The one point between them is not resolved either ($p = 0.73$, Figure 1; a
95\% deficit of up to 4.8 points allowed). The costs are calculated with a rate of \$0.35/\$2.75 per MTok from a third-party OpenRouter host \cite{openrouter2026}. Our test was run locally, and anyone with capable hardware could potentially host it with a lower cost.

Third, the two GPT-5.6 models provide the clearest case in the paper for
spending on the scaffold rather than on the model. GPT-5.6-Luna with a manager
matches GPT-5.6-Terra's single call on accuracy at 44\% of the price: $77.8$ against
$77.0$, for \$1.50 against \$3.41. The price half is decisive - \$1.91 a pass,
$p = 2.1\times10^{-7}$ across runs and $p < 3.5\times10^{-5}$ across problems (Table 4). The 0.8-point lead is not
itself significant (two-sided $p = 0.76$), but there is no sign of a deficit either, and the one-sided 95\% bound rules out Luna's manager trailing by
more than 2.4 points.

\FloatBarrier

\subsection{The effect of truncation}

An empty solution is an automatic fail, so any condition that emits fewer empty solutions
gains pass@1. Across the seven pinned-backend arms of \S2.1 we can separate that channel from
genuine problem-solving exactly, because these runs carry per-call
\texttt{finish\_reason}. Two counts are involved:
a cap hit is a graded generation that reached the token limit
(\texttt{finish\_reason=length}) - for the single arms that is their one call, for the
manager its final answer call - while no code is a problem-pass that ended with
nothing the extractor could grade. A call can hit the cap and still be graded, because the
model often writes a complete solution inside its reasoning before the cap arrives.
Out of 500 problem-passes per arm, single (s) / manager (m):

\begin{table}[htbp]
\centering
\caption{\textbf{Cap hits and empty solutions per arm.} Single (s) / manager (m), out of
500 problem-passes each.}
\small
\begin{tabular}{@{}lcccc@{}}
\toprule
\textbf{Model} & \textbf{Cap hits s / m} & \textbf{No code s / m} &
\textbf{of which refusals} & \textbf{Cap} \\
\midrule
GPT-5.6-Terra  & 0 / 0            & 0 / 0          & 0 & 128k\textsuperscript{b} \\
GPT-5.6-Luna   & 0 / 0            & 0 / 0          & 0 & 128k\textsuperscript{b} \\
Qwen3.8-27B    & \textbf{150} / 5 & \textbf{35} / 0 & 0 & 128k\textsuperscript{a} \\
Claude Fable 5 & 3 / -            & 9 / -          & \textbf{6} & 128k\textsuperscript{b} \\
\bottomrule
\end{tabular}

\vspace{0.6ex}
\begin{minipage}{0.92\linewidth}
\footnotesize
\textsuperscript{a}~Every count in this table is at a 128k cap, so the rows are
like-for-like. Qwen3.8's manager arm ran natively at 128k; its single arm was run at 250k
and is counted here cap-matched back to 128k. The 250k counts are given below.\\
\textsuperscript{b}~A \textit{hard} ceiling rather than a chosen setting: the GPT-5.6
family and Fable 5 both cap a synchronous API response at 128k output tokens (\S3.2), so
for these three arms a truncation cannot be relieved by raising the cap. Qwen3.8-27B, served
locally, has no such ceiling - which is why its single arm could be run at 250k at all.
\end{minipage}
\end{table}

\textbf{The two OpenAI arms never truncate and never emit an empty solution} - zero cap
hits and zero empty answers across all 2{,}000 problem-passes. Their $\Delta$s in \S2.1
therefore contain no rescue component at all.

\textbf{On Qwen3.8-27B the two counts differ by more than a factor of four.} At 128k its single arm is
cut off on 150 of 500 problem-passes, yet only 35 end with nothing to
grade. The large majority of cut-off generations still contain a complete solution, written
inside the reasoning stream before the cap arrived and recovered from there by the
extractor. A cut-off rate is not a loss rate. Against those 35 the manager has zero: it is not truncation-free
internally - 108 of its 3{,}235 calls across the five passes hit the cap, 95 of them
workers, touching 75 of the 500 problem-passes - but the loop
absorbs them, because a cut-off worker costs a round, not the answer.

\paragraph{Exploratory: the same run with the 128k limit lifted.}
Qwen3.8-27B was run locally at its native context window of 262{,}144 tokens \cite{qwen2026c}, so it was able to generate more than 128k tokens. Its single arm was in fact generated at 250k, with a 128k output cap imposed after the fact. Reading these outputs to their natural end recovers what the limit was hiding. We treat the reading as
exploratory - it is cap-matched to nothing else in the paper, so it cannot carry a
$\Delta$ against the 128k manager arm - but it bounds what the extra budget is worth:

\begin{table}[htbp]
\centering
\caption{\textbf{Qwen3.8-27B's single arm read at both caps.} Cap-matched to 128k, and as
generated at 250k.}
\small
\begin{tabular}{@{}lccc@{}}
\toprule
\textbf{Qwen3.8-27B single} & \textbf{pass@1} & \textbf{Cut off by the cap} &
\textbf{No code at all} \\
\midrule
128k (cap-matched, \S2.1)       & $63.0 \pm 4.1$ & 150 / 500 \ (30.0\,\%) & 35 / 500 \ (7.0\,\%) \\
250k (as generated)             & $65.6 \pm 4.6$ & 124 / 500 \ (24.8\,\%) & 21 / 500 \ (4.2\,\%) \\
\bottomrule
\end{tabular}
\end{table}

The tokens past 128k are worth 2.6 points, and they earn them by shrinking the
tail rather than by lifting the run as a whole: 26 fewer generations end mid-stream and 14
fewer problem-passes finish with nothing to grade, with the per-pass no-code count falling
5/5/6/7/12 $\to$ 4/2/3/4/8. The pass that keeps the worst tail is also, at 128k, the joint
lowest-scoring of the five. So on this model the cap acts less like a score multiplier than like a tax on
the hardest few problems per pass.

\paragraph{Manager ``rescue,'' measured directly.}
Of the single Qwen3.8-27B arm's 35 no-code cells, the manager passed 25, failed 10, and left
none empty. This mechanism results in a $25/500=5.0$ point increase to the manager's score, about a fifth of the manager's +23.4 point total improvement over the single pass.


\paragraph{What the empty cells cost each model.}
Re-scoring each single arm over only the problem-passes where it emitted code measures the
same effect for every row:

\begin{table}[htbp]
\centering
\caption{\textbf{Each single arm re-scored} over only the
problem-passes where it emitted code.}
\small
\begin{tabular}{@{}lccc@{}}
\toprule
\textbf{Model (single arm)} & \textbf{As scored} & \textbf{Emitted code} &
\textbf{Restricted to those} \\
\midrule
GPT-5.6-Terra  & 77.0 & 500/500 & 77.0 \\
GPT-5.6-Luna   & 67.2 & 500/500 & 67.2 \\
Qwen3.8-27B    & 63.0 & 465/500 & 67.7 ($+4.7$) \\
Claude Fable 5 & 87.4 & 491/500 & 89.0 ($+1.6$) \\
\bottomrule
\end{tabular}
\end{table}

\paragraph{Refusals - a failure mode unique to Fable 5.}
Six of Fable 5's nine empty cells are not truncations but safety refusals: the
API returned \texttt{stop\_reason=refusal} with no content. All six fall on three
ordinary competitive-programming problems, none of them security- or biology-adjacent, and
none is deterministic - problem 3739 was refused in three of five passes, 3682 in two,
abc393\_e in one. We score them as failures and do not configure a fallback model,
because a cell labelled Fable 5 has to contain Fable 5's outcome; a silent rescue by a
different model would contaminate the measurement. The cost is about 1.2 points of the
$87.4$, and it is a penalty no other arm in this paper pays. The remaining three empty
cells are true truncations, all on a single problem (arc191\_d) that spends the full 128k
reasoning without emitting an answer - the cap bounds thinking and answer together, so
that is a possible outcome rather than a malformed
response \cite{anthropic2026effort}.

\FloatBarrier

\subsection{Earlier results - the OpenRouter-served model set}

The five models below were served through the OpenRouter gateway and run on the
original scaffold, so they are not directly comparable to the pinned-backend arms of
\S2.1: both the serving path and the scaffold version differ (\S3).

\begin{table}[htbp]
\centering
\caption{\textbf{LCB-100 pass@1 (\%), single $\to$ manager.} For the OpenRouter-served
models on the original scaffold. ``-'' = not run; Qwen3.5-9B thinking-on returns
reasoning-only replies and is unusable. The seven pinned-backend, five-pass arms are reported separately in
\S2.1 and are not pooled here, because both the serving path and the scaffold version
differ.}
\small
\begin{tabular}{@{}llccc@{}}
\toprule
\textbf{Model} & \textbf{Params} & \textbf{128k $\cdot$ ON (1 pass)} &
\textbf{128k $\cdot$ OFF (1 pass)} & \textbf{16k $\cdot$ OFF ($\times$5)} \\
\midrule
Opus-5        & n/a            & 85 $\to$ 91 (\textbf{+6})            & -                       & - \\
Kimi-K3       & ${\sim}2.8$T   & 83 $\to$ 82 ($-1$)                   & 32 $\to$ 74 (\textbf{+42}) & 32.2 $\to$ 62.6 (\textbf{+30.4}) \\
Minimax-M3    & 428B           & 60 $\to$ 66 (+6)                     & 25 $\to$ 37 (\textbf{+12}) & 21.2 $\to$ 32.2 (\textbf{+11.0}) \\
Qwen3.6-35B   & 35B            & 25 $\to$ 43 (\textbf{+18})           & 35 $\to$ 26 ($-9$)         & 27.8 $\to$ 26.6 ($-1.2$) \\
Qwen3.5-9B    & 9B             & \textit{unusable}                    & 17 $\to$ 20 (+3)           & 14.6 $\to$ 21.8 (\textbf{+7.2}) \\
\bottomrule
\end{tabular}
\end{table}

\input{fig-128k-reason-on}

With thinking on:

\begin{enumerate}
  \item \textbf{Opus-5 tops the study at 85 $\to$ 91.} The $+6$ is a net of seven problems
  gained and one lost.

  \item \textbf{The largest gains come from smaller models.} Ordered by scale, the single
  curve rises $+60$ pts (35B$\to$Opus) and the manager curve $+48$ pts (Figure 4).
\end{enumerate}

% Flush Figure 4 before the thinking-off material queues Figures 5 and 6 behind it. Without
% this the three compete for float pages in order and each lands a page or two after the
% text that cites it.
\FloatBarrier

\paragraph{Thinking off - cheaper runs, more passes.}
Reasoning-off runs are far cheaper (no long reasoning tokens), so we ran
five independent passes per condition at a 16k cap, plus a one-pass 128k.

\paragraph{Significance.}
Run-to-run spread across the five 16k passes is small - per-condition SD of 0.5--4.8 points
single, 1.8--3.6 with the manager - so the deltas are not noise. Per model we run a
paired sign-flip permutation test with the problem as the unit ($n = 100$),
Holm-corrected across the four models, and report 95\% confidence intervals across the five
passes (Figure 5). The manager effect is significant for three of four models:

\begin{table}[htbp]
\centering
\caption{\textbf{Manager $-$ single over the five 16k passes, reasoning off.}}
\small
\begin{tabular}{@{}lccc@{}}
\toprule
\textbf{Model} & \textbf{$\Delta$ (mgr $-$ single)} & \textbf{Holm $p$} & \textbf{Significance} \\
\midrule
Kimi-K3     & $+30.4$ & $< 2\times10^{-5}$ & *** \\
Minimax-M3  & $+11.0$ & $6\times10^{-5}$   & *** \\
Qwen3.5-9B  & $+7.2$  & $4\times10^{-4}$   & *** \\
Qwen3.6-35B & $-1.2$  & $0.7$              & n.s. \\
\bottomrule
\end{tabular}
\end{table}

\paragraph{Differences between models.}
A Tukey HSD separates Kimi-K3 from all three others, and Minimax-M3 from Qwen3.6-35B
($p = 0.0038$); Qwen3.5-9B sits between them and separates from neither ($p = 0.71$,
$0.087$). So the manager helps Kimi most and Minimax-M3 more than Qwen3.6-35B; where the
9B belongs cannot be resolved from 100 problems.

% Each fig-*.tex is a [p] float, so it cannot be placed before the point it is \input.
% These sit directly under the paragraph that cites them: further down the source and the
% float queue backs up, which is how Figure 5 ended up three pages after its reference.
\input{fig-16k-reason-off}

\paragraph{128k control.}
One reasoning-off pass at the larger cap largely eliminates truncation and reproduces the
ordering - Kimi $+42$, Minimax $+12$, Qwen3.5-9B $+3$, Qwen3.6-35B $-9$ (Figure 6).
Empty output is not fully removed: single-arm emit rates run 79--96\% against 92--100\%
for the manager, so part of the gap at the small end remains an emit-rate effect. It
cannot explain Kimi, where a 4-point emit gap accompanies a $+42$-point score gap. Unlike
the thinking-on setting, the single and manager curves diverge here (single $+15$ pts,
manager $+54$ pts end to end), though the two settings span different model sets. From 35B
upward the manager's advantage grows with model strength rather than converging.

\input{fig-128k-reason-off}

\paragraph{Truncation in this model set.}
The 16k runs predate the per-call \texttt{finish\_reason} instrumentation of \S2.3, so the
truncation channel can only be bounded here, not decomposed. At 128k it is negligible
everywhere
except Qwen3.6-35B with thinking on, which truncates on 48 of 100 single-call
problems because its reasoning is clamped at a provider ceiling of 32k, far below the cap -
an artifact of the gateway, not of the model. At the tight 16k cap the single arm produces
no code on 13--43 problems depending on the model while the manager loses far fewer
(Qwen3.5-9B $43 \to 7$, Minimax $35 \to 5$, Kimi $13 \to 1$). And the problems lost to the
cap are the hard ones: those producing no code in at least three of the five 16k passes go
on to pass at 0--25\% at 128k, against 27--37\% for the rest, so the manager's emit-rate
advantage is concentrated where the answers were hardest to reach anyway.

Rescue is directly observable on the one model with enough truncation to measure it: of
Qwen3.6-35B's 48 single-truncated problems the manager passed 21, left 15 still empty and
filled but failed 12. Opus-5 truncated 3 and the manager passed all 3; Minimax-M3 and
Kimi-K3 never truncate at 128k, so for them there is nothing to rescue and their deltas
cannot be attributed to this channel at all.

\section{Method}

\S3.1 describes the current scaffold, v2, which produced the pinned-backend results
in \S2.1 and is shown in Figure 7. An
original version produced the OpenRouter-served set of \S2.4. It is the same loop
and the same prompts, with four things absent, all of them on the manager arm:

\begin{center}
\small
\begin{tabular}{@{}lll@{}}
\toprule
& \textbf{v2 (\S3.1, \S2.1)} & \textbf{original (\S2.4)} \\
\midrule
Round budget (\texttt{MAX\_ITERS}) & 10 manager$\to$worker cycles & 4 \\
Sample-test verifier (step 5)      & yes                          & absent \\
Cut-off summarizer                 & yes                          & absent \\
Workspace files feeding a prompt   & size-bounded                 & unbounded \\
\bottomrule
\end{tabular}
\end{center}

Because the single-call baseline is one call under either version, the difference touches
only the manager arm - so manager-minus-single deltas are not strictly comparable across
the two groups, and we report them as separate conditions rather than pooling them into
one table.

\subsection{The manager--worker loop}

Every role uses the same underlying model, called in a fresh context and coordinating only
through the shared workspace:

\begin{center}
\small
\begin{tabular}{@{}ll@{}}
\texttt{<ws>/task.md}    & the problem statement \\
\texttt{<ws>/plan.md}    & the manager's overarching plan \\
\texttt{<ws>/tasks.json} & the task list \texttt{[\{id, desc, status, result\}]} \\
\texttt{<ws>/notes.md}   & accumulated ideas / findings / partial proofs \\
\texttt{<ws>/solution.py} & current best code \\
\end{tabular}
\end{center}

The control flow (Figure 7) is:

\begin{enumerate}
  \item \textbf{Manager - plan.} The manager reads the problem and writes a 3--6
  sentence strategy plus 3--6 concrete seed tasks.

  \item \textbf{Worker - brainstorm (ideation).} The first worker does \textit{not}
  write a solution; it identifies the core difficulty, lists candidate approaches and
  pitfalls, and appends them to \texttt{notes.md}, proposing next steps.

  \item \textbf{Manager - manage (loop).} The manager folds the plan and the brainstorm
  into one curated task list (merging duplicates, marking done items, adding only
  genuinely new sub-tasks), then either declares the problem done or names the
  single next task.

  \item \textbf{Worker - do the task.} A fresh worker executes that one task, rewrites
  \texttt{solution.py}, appends what it did to \texttt{notes.md}, and proposes remaining
  steps.

  \item \textbf{Verifier - run the sample tests} \textit{(absent from the original
  scaffold, \S3)}. Whenever the round's worker produced a fresh candidate,
  the program is executed against the problem's public sample tests - the stdin-format
  tests, which cover 73 of the pinned 100 problems; the 27 LeetCode-style problems carry
  functional/call-based public tests the engine does not execute, so they are checked only
  by the hidden-test grader. The pass/fail verdict, with the first failing case, is fed
  back to the manager and treated as ground truth: a failing run overrides a
  done verdict, forcing the loop to continue with a fix-or-switch task. Control
  returns to the manager (step 3).

  \item \textbf{Finalizer.} A finalization worker emits the definitive solution whenever
  the loop ends without a clean sign-off - the round budget is spent, the manager
  reissues a task it just handed out, or it names no task at all. The call is skipped when
  the manager declares the problem done and a usable solution is already on disk, so a
  redundant final pass cannot overwrite a correct answer.
\end{enumerate}

% Single portrait panel (aspect 0.75), so it is bounded by height like the two-panel
% figures above: 0.88\textheight is 7.9in tall and 5.9in wide, leaving room for the
% caption on a float page.
\begin{figure}[p]
\centering
\includegraphics[width=\linewidth,height=0.88\textheight,keepaspectratio]%
  {\plotdir/how_the_manager_arm_answers_one_question_multiagent_py_current_tree_light.png}
\caption{\textbf{Manager--worker control flow}, as the v2 scaffold runs it. Every box is
one model call in a fresh context except the sample-test run, which executes the candidate
program; the workspace files are the only state any role sees. The original scaffold is
the same loop with the round budget at 4 and without the two steps v2 adds: the
sample-test verifier (step 5) and the cut-off summarizer.
(\texttt{plot\_agent\_loop\_flowchart.py}.)}
\end{figure}

Guards keep the loop cheap and safe: a round budget of 10 manager$\to$worker
cycles (\texttt{MAX\_ITERS}); a no-progress guard (if the manager reissues the
exact task it just handed out, the loop stops); and a cut-off summarizer - a
worker that hits the token cap mid-attempt has its partial thinking summarized by a fresh
short call so its ideas still reach the manager. The original scaffold ran the same
no-progress guard, but with the budget at 4 and no summarizer (\S3).

The single-call baseline is the same model, at the same temperature as the
manager arm's workers (0.2), given the same problem in exactly one call - no shared
workspace, no loop, and no other role in the prompt. The manager's generative calls 
use slightly higher temperatures, where the work is generative rather than executive: 0.3 to write the plan and 0.4
to brainstorm, against 0.2 for task execution and for curating the task list. These
settings are identical in both scaffold versions. It receives the solver system prompt
verbatim; the manager arm's workers receive that same prompt wrapped in a subagent
preamble and a four-section output contract, and their user message additionally carries
the plan, the accumulated notes and the current artifact. The two conditions differ only
in that scaffold.

% Flush the flowchart before the protocol section. It is declared above as a [p] float and
% LaTeX was carrying it forward onto the next page, which landed it in the middle of the
% Conditions paragraph and split that paragraph off from its own itemized list.
\FloatBarrier

\subsection{Benchmark and protocol}

\paragraph{Benchmark:}
LiveCodeBench code-generation, release\_v6 \cite{jain2024}. We take
the 100 latest problems in the hard split (by contest date). Using the latest problems
reduces training-data contamination \cite{jain2024}. Solutions are graded by
LiveCodeBench's own hidden-test evaluator.

\paragraph{Models:}
Qwen3.5-9B \cite{qwen2026a} (9B), Qwen3.6-35B-A3B \cite{qwen2026b} (35B total, 3B active
per token), Qwen3.8-27B \cite{qwen2026c} (27B, open weights, served locally in FP8), Minimax-M3
\cite{minimax2026} (428B total, ${\sim}23$B active), Kimi-K3
\cite{moonshot2026} (${\sim}2.8$T total, ${\sim}16$ of 896 experts active per token), and
the closed frontier models Opus-5 \cite{anthropic2026} (size undisclosed), two GPT-5.6
variants, Terra and Luna \cite{openai2026} (sizes undisclosed), and Claude Fable 5
\cite{anthropic2026} (size undisclosed) - nine model configurations in all.

Serving splits the set in two, separating models in \S2.1 from
\S2.4. The five models of \S2.4 were served through a single OpenAI-compatible gateway
(OpenRouter) on the original scaffold, so that within that set the only thing varying is
the model. The four models of \S2.1 are each served on a pinned backend under the v2
scaffold (\S3) - Terra and Luna by the OpenAI API, Qwen3.8-27B by our own vLLM, Fable 5 by
the Anthropic Messages API - because gateway routing proved to be the dominant source of
run-to-run noise (\S4.5). Fable 5 was run single-only, so it contributes a single-call
figure and no manager delta.

\paragraph{Conditions:}
The pinned-backend set of \S2.1 runs at one setting - a 128k output cap, reasoning on,
$\times5$ passes. Terra, Luna and Qwen3.8-27B are run in both
single and manager conditions; Fable 5 is run single-only, as a reference arm. The cap is a
per-call \texttt{max\_tokens} bound on generation, not a context-window limit. 128k is
also the ceiling some providers impose on a single response - Opus-5, Fable 5 and the
GPT-5.6 family all cap a synchronous API response at 128k output tokens
\cite{anthropic2026,openai2026}, so for those arms the cap is the provider's hard limit
rather than an experimental choice. Qwen3.8-27B, served locally, has no such response
ceiling; what bounds it instead is its 262{,}144-token context window
\cite{qwen2026c}, which prompt and output share. The 250k output setting used for \S2.3 is close to the largest output the model can produce at
all. The gateway-served set of \S2.4 are run at the three settings below:

\begin{itemize}
  \item 128k output cap, reasoning on, $\times1$ pass - thinking enabled (native
  \texttt{reasoning\_effort} for models that support it; a bounded 20k reasoning budget
  for open models that do not).

  \item 16k output cap, reasoning off, $\times5$ passes - 5 independent passes per
  condition to measure run-to-run variance and support paired significance tests.

  \item 128k output cap, reasoning off, $\times1$ pass - removes the small cap as a
  confound.
\end{itemize}

\paragraph{Thinking:}
The output cap and the scaffold are matched across the seven pinned-backend arms of \S2.1.
Thinking depth is not, because no two of these providers expose the same control
over it. Qwen3.8-27B, on our own vLLM, is sent no \texttt{reasoning\_effort} at all: the
Qwen chat template leaves thinking on and unbudgeted, so the only thing bounding it is the
128k \texttt{max\_tokens}, which on this stack covers reasoning and answer together. The two
GPT-5.6 arms are likewise sent no \texttt{reasoning\_effort}, leaving them at OpenAI's model
default. Fable 5 has
\texttt{thinking} set to \texttt{adaptive} with \texttt{display: summarized} - the most this API
hands back, since no setting returns the raw chain of thought - and
\texttt{output\_config.effort} set as \texttt{high}, which is also the API
default \cite{anthropic2026effort}.

\paragraph{Instrumentation:}
In the two 128k-cap conditions, every model call's \texttt{finish\_reason} and
completion-token count is logged, and each problem's final record carries a status in
\{\texttt{ok}, \texttt{truncated}, \texttt{empty\_stop}, \texttt{empty}, \texttt{error}\}.
The 16k $\times$ 5-pass runs predate that instrumentation: their records carry only the
extracted code and its pass/fail, so per-call outcomes cannot be recovered for that
condition. An empty code field is considered a failure.

\paragraph{Cap-matching:}
One arm is not natively cap-matched to its partner: Qwen3.8-27B's single call was generated
at 250k while its manager arm ran at 128k (\S2.1). To compare them at one cap we replay the
stored generations rather than re-running the model. Each generation is truncated to
128,000 output tokens and the solution re-extracted from that prefix, after which it is
scored on the same evaluator as every other cell (\S3.3). Three details matter for
interpreting the result:

\begin{itemize}
  \item \textbf{Truncation is token-exact}, using the serving stack's own tokenizer, not a
  character-count approximation. The cases that decide the outcome are exactly those where
  a complete code block sits just before or just after the boundary, so an approximate cut
  would misclassify precisely the problems the procedure exists to settle.

  \item \textbf{What is truncated is the whole output stream, reasoning then answer},
  because that is what the original cap bounded. Where the cut lands decides what the
  extractor sees: past the reasoning it sees a truncated answer; inside the reasoning the
  answer does not exist at all, and the harness's empty-content fallback hands it the
  truncated reasoning instead. That fallback is why so many cut-off generations still yield
  gradeable code (\S2.3).

  \item \textbf{The cap is invisible to the model, which is what makes the replay exact.}
  This arm ran on our own vLLM, where \texttt{max\_tokens} is a parameter the
  server enforces: it never enters the prompt and does not affect the model's output. Generation simply cuts off when the limit
  is reached. The records show that 124 of the 500 calls end at
  \texttt{finish\_reason=length}, cut off abruptly rather than wound up. Nor is this local to vLLM: surfacing a budget to a model takes a
  separate, opt-in mechanism, and Anthropic ships \texttt{task\_budget} precisely so a model
  can ``finish gracefully \dots rather than cutting off mid-action,'' while
  \texttt{max\_tokens} remains the enforced ceiling that ``truncates the
  response'' \cite{anthropic2026effort}. There is therefore no difference from genuine re-run at 128k.
\end{itemize}

\paragraph{Statistics:}
Reported $\pm$SD is the sample SD (divisor $n - 1 = 4$) of the five pass scores. Five-pass
comparisons: paired sign-flip permutation test on per-problem mean $\Delta$ ($n = 100$
problems; $2\times10^{5}$ resamples), Holm-corrected where a family of models is tested;
95\% t-intervals for the mean pass@1 across the five independent passes ($\mathrm{df} = 4$). Pooled problem$\times$pass agreement tables use exact
McNemar on the pooled discordants.

\subsection{A correction to the evaluator}

While inspecting transcripts for \S4.1 we found a defect in the LiveCodeBench harness
itself, and every number in \S2.1--\S2.3 is reported after fixing it. The evaluator does
not run a candidate as a subprocess; it executes it in-process with \texttt{sys.stdin}
replaced by a mock. That mock's binary view implemented \texttt{readline()} as
\texttt{inputs.split(b"\textbackslash n")[0]} - a \textit{stateless} expression that
returns the first line on every call. A program reading multi-line input through
\texttt{sys.stdin.buffer.readline()} therefore read line 1 repeatedly and scored wrong no
matter how correct it was. The text-mode \texttt{sys.stdin.readline} was patched to a
proper iterator and behaved correctly, and \texttt{buffer.read()} returns the whole payload
and was also unaffected, which is why the common
\texttt{sys.stdin.buffer.read().split()} idiom never exposed it. We replaced the mock's
binary view with one backed by a \texttt{BytesIO} so that reads advance a position, and
re-scored every stored generation; no model was re-run.

This defect warrants explicit reporting for two reasons. First, it interacts
badly with the v2 verifier: that step runs the candidate as a \textit{real} subprocess
(\S3.1, step 5), where \texttt{buffer.readline()} works correctly, so the manager was told
its solution passed the public samples for programs the grader then marked wrong - the one
signal in the loop meant to be external and trustworthy, certifying the wrong answer.
Second, the exposure is uneven across models rather than a constant offset: 311 of the
3{,}456 \S2.1 outputs containing code use the idiom, and 97\% of those were scored wrong
against 20\% for the outputs that do not. Fable 5 never uses it and its scores are unchanged to the
decimal. The \S2.4 model set reaches for it once in 5{,}012 generations, and then through
an alias rather than the literal spelling - Kimi-K3's single call on abc396\_g binds
\texttt{sys.stdin.buffer} to a name and calls \texttt{readline()} on that. Every \S2.4
number is re-scored on the fixed evaluator all the same, and exactly one cell moves: Kimi's
128k thinking-on single arm, from 82 to 83. Because usage of the idiom is a property of a model's coding style, a
harness bug of this shape is not a wash across a leaderboard - it silently penalises the
models whose style happens to trip it.

\section{Discussion}

\subsection{Where the manager clearly helps (transcript evidence)}

Each example is a problem that the single call failed and the manager solved. We include
one per model, with each illustrating a different mechanism.

\begin{itemize}
  \item \textbf{Surfacing the efficiency trap for a weak model} \textit{(Qwen3.5-9B, 128k,
  reasoning off, LCB abc385\_d).} On a path-simulation problem that counts houses lying on
  each move segment, the single 9B pass \textit{timed out} - it ran a na\"ive
  per-segment scan. The manager's brainstorm flagged, in writing, that this check is
  $O(N\cdot M) \approx 4\times10^{10}$ and prescribed a range-query structure
  \textit{before} any code, which the workers then filled in. For a small model, the
  workspace supplies exactly the up-front planning a single pass skips (cf.\ \S4.3).

  \item \textbf{Escaping a stuck / over-long single pass} \textit{(Qwen3.6-35B, reasoning
  on, LCB abc394\_f).} The single call was cut off mid-reasoning at 32,768 tokens and
  returned nothing at all. The manager solved it in one worker cycle: the brainstorm
  crystallized the structural insight (an ``alkane'' subtree needs a degree-4 centre and
  $\ge 5$ vertices) and the worker implemented an iterative DFS to avoid the
  recursion-depth failure that also affects many single calls. Across all manager-arm calls for this model,
  the same 32k ceiling was hit 117 times. With several
  attempts per problem, one clamped call no longer costs the problem.

  \item \textbf{Splitting a two-objective algorithm into separate sub-DPs}
  \textit{(Minimax-M3, reasoning on, LCB 3701).} The task needs both the minimum edit cost
  of a ``good caption'' (runs of equal letters, each $\ge 3$ long) and, among optimal
  captions, the lexicographically smallest. The single-pass program printed nothing on a
  small case (\texttt{cdcd}, whose answer is \texttt{cccc}) - its
  lexicographic-reconstruction branch was broken. The manager split the coupled objectives
  across worker cycles - one worker implementing the forward cost DP over capped
  run-lengths (\texttt{dp[i][c][k]}), a separate worker a suffix DP used purely for the
  lexicographic reconstruction - turning one tangled solution into two
  independently-written pieces.

  \item \textbf{Forcing an explicit reduction before coding} \textit{(Kimi-K3, 128k,
  reasoning off, LCB 3687).} The single pass returned a wrong answer that
  \textit{over-counted} the longest unique-value path (9 and 3 where the answers were 6
  and 2), never shrinking its window on a repeated value. The manager's \textit{brainstorm}
  step first wrote the reduction into \texttt{notes.md} - ``this is
  longest-substring-without-repeating-characters along each root-to-leaf path; maintain a
  window start via last-seen depths; $O(n)$ DFS'' - and only then did a worker implement
  it, over two manager$\to$worker cycles. The scaffold turned an implicit leap into an
  explicit, written plan the next call could build on.

  \item \textbf{Replacing an over-engineered structure with a smaller one}
  \textit{(GPT-5.6-Terra, reasoning on, LCB 3688).} The task is a maximum-subarray sum
  where all occurrences of one chosen value may first be deleted. The single pass wrote
  4{,}615 characters implementing a segment tree carrying minimum, second-minimum, maximum
  and a lazy add - a Segment-Tree-Beats-shaped structure indexed by candidate negative
  value - and got it wrong. The manager's plan named the actual difficulty in advance: the
  state must be ``compressed so it does not require maintaining all distinct values per
  index''. Nine calls later the manager arm returned 1{,}860 characters: one
  segment tree over the standard (total, best-prefix, best-suffix, best-subarray) merge.
  The scaffold's contribution here is \textit{subtraction} - a correct solution 2.5$\times$
  smaller than the failed one - which is the opposite of the usual worry that a scaffold
  adds machinery. Manager-only win in four of five passes.

  \item \textbf{Deriving the bounding lemma before coding} \textit{(GPT-5.6-Luna, reasoning
  on, LCB abc397\_e).} A tree on $NK$ vertices must be decomposed into $N$ paths of exactly
  $K$ vertices. The single pass carried a \textit{set} of unfinished path lengths per
  subtree (\texttt{open\_sets}) - a general state it never managed to close - across 4{,}092
  characters. The manager's brainstorm first established the lemma that makes the problem
  tractable: at any vertex there are at most two unfinished paths, and they must
  either extend through the vertex or join into one complete path of exactly $K$. With the
  branching bounded in writing, a worker implemented the postorder scan directly, in 2{,}166
  characters over five calls. Manager-only win in two of five passes.

  \item \textbf{Rescuing a call that dissolved into self-verification}
  \textit{(Qwen3.8-27B, reasoning on, LCB abc397\_g).} Maximise the shortest 1-to-$N$
  distance after raising exactly $K$ of $M$ edges to weight 1. The single call spent all
  250{,}000 tokens and never emitted a program: what the extractor recovered is the tail of
  its own edge-case interrogation (``\textit{Now, let's think about if the min-cut graph's
  min cut capacity is not equal to min cost labeling \dots{} We proved. Good.}''), the
  \S4.2 failure mode in its pure form. Strikingly, that reasoning had \textit{already
  reached the right idea} - it is discussing the min-cut formulation - and could not stop to
  write it down. The manager arm committed the same reduction to \texttt{plan.md} as a
  plan (binary-search the target distance $D$; feasibility is a vertex-labelling minimised
  by an $s$-$t$ cut on a layered graph), then had a worker implement it: five calls,
  3{,}548 characters, and a manager-only win in all five passes. The scaffold's
  role is not supplying the insight but forcing it onto disk before the budget runs out.

  \item \textbf{Writing the proof before the code} \textit{(Opus-5, reasoning on, LCB
  abc388\_e).} On maximising the number of disjoint (top, bottom) pairs with
  $2\cdot\text{top} \le \text{bottom}$, the single pass emitted code whose feasibility test
  \textit{over-counted} - reporting 225 pairs where 220 is optimal. The manager first
  recorded the \textit{justification} in \texttt{notes.md} - an exchange argument that
  the $K$ smallest elements as tops, matched in sorted order to the $K$ largest as
  bottoms, is optimal, and that feasibility is monotone in $K$ (licensing a binary search)
  - then had one worker implement it and a second harden it. This is a genuine
  correction, not truncation relief: separating the argument from the implementation
  caught an error the one-shot attempt made.
\end{itemize}

\subsection{Truncation and large-context management}

With multiple workers, organized notes replace an overflowing context window. State lives
in \texttt{plan.md}, \texttt{tasks.json}, and \texttt{notes.md} rather than in one
ever-growing transcript. Each worker sees a compact, curated view (plan + notes + current
solution + one task) instead of the full history, which both lowers per-call length and
keeps the salient facts in view. The workspace is, in effect, external memory that the
token cap cannot truncate.

A crowded context or an unbounded generation does not merely risk being cut off - it can
actively \textit{degrade} a model, and the effect may be worse in the smaller models. Left
to fill a long output, a weaker model risks losing the thread and looping. On
\texttt{abc399\_e} (Qwen3.8-27B single, reasoning on, 250k cap) the model produced 675{,}000
characters of reasoning, spent the entire 250{,}000-token budget, and emitted no solution
at all; inside that stream it repeated the single line \textit{``Alphabet 5 with
a$\leftrightarrow$b, c$\leftrightarrow$d, e$\to$e fixed, and one letter? no.''}
7{,}743 times. Eight of the fourteen generations that spent the full budget without
emitting code loop on the same template - \textit{``Now, let's consider if there is a
possibility of \dots{} no.''} - inventing hypothetical edge cases against a candidate
solution and answering each one, hundreds to thousands of times over, without ever
terminating. This is the mechanism behind the otherwise odd finding of \S2.3 that most
cut-off generations still contain a complete solution: the model had reached an answer and
then talked past the cap re-checking it. The manager breaks the work into short,
self-contained calls, each far from the cap, so the model's work actually reaches the disk. Bounding each
call's workload, context and length is thus a guard against this collapse as well as a cost
saving.

The same mechanism extends beyond this benchmark. A task whose inputs, working state, or
intended output simply do not fit - a codebase larger than the context window, or a
deliverable longer than the per-call output ceiling described in \S3.2 - is not solvable
in one call at any reasoning effort. Decomposing the task into a sequence of bounded steps
allows the model to complete more complex tasks that would exceed a model's context window
or API thinking/output limit outright.

\subsection{A large effect on thinking-disabled models}

The manager provides the \textit{largest} benefit when the base model is least effective at
self-organizing its own reasoning, and that weakness takes two forms. The first is having
no internal planning stage at all: with thinking disabled, a single call jumps straight to
code, and the scaffold's brainstorm-then-plan structure gives the model an
alternative space to plan - on disk, across calls - that it would otherwise
not use (Kimi $+42$, Minimax $+12$ at 128k-off). The second is thinking that runs away:
Qwen3.6-35B with reasoning on spends its whole budget deliberating - 48 of its 100
single calls were cut off mid-thought, each at a provider-side 32,768-token output clamp,
which is therefore that arm's modal completion length by a wide margin - and it gains $+18$, its
largest gain in any condition. Where a model can both plan internally \textit{and} stop on
its own, the scaffold has less to add. The workspace
substitutes for the reasoning a model cannot organize for itself.

\subsection{Regressions: when the manager hurts}

The manager is not free. Its deliberation can settle on a worse answer than a single call
would have produced:

\begin{itemize}
  \item \textbf{Deliberating its way into an algorithm it had already rejected.} On LCB
  3765 (Qwen3.6-35B, 128k-off) the single pass wrote a correct convex-hull-optimized
  $O(n^2)$ dynamic program. The scaffold's ideation stage identified that same
  optimization and talked itself out of it - ``implementing CHT is complex and
  error-prone'' - committing instead to an $O(n^3)$ table its own notes call ``too slow
  for Python,'' on the reasoning that the ``test cases are weak.'' The one worker round
  then implemented that plan with a further bug, returning \texttt{dp[n][n]} - the
  all-singletons partition - rather than the minimum over subarray counts its own task
  list specified. The result is both wrong and, at 19\,s for $n = 1000$ against 0.3\,s,
  far too slow. Qwen3.6-35B is the standout loser with reasoning off ($\Delta$ of $-1.2$
  at 16k and $-9$ at 128k): the compact, correct code is more often \textit{disturbed}
  than helped by decomposition.
\end{itemize}

In practice, the scaffold can backfire when deliberation produces a worse plan than the
model's initial approach and the manager fails to detect the regression.

\textbf{Two inexpensive, training-free additions may address these failures}, one for each part of the
failure above - the plan that was wrong, and the implementation that went unchecked.
First, \textit{fresh-perspective workers}. Because every worker inherits the accumulated
notes and the current solution, a wrong early approach anchors everything downstream -
the failing rewrite of 3765 built on an approach its own ideation stage had already judged
too slow, instead of reconsidering it. Spawning some workers with the raw problem
statement and no prior context would give the manager an independent attempt to
compare against the evolving one, and keep whichever is better; in effect this protects
the tight single-shot solution that the scaffold otherwise disturbs, and diversifies away
from a bad initial framing (the intuition behind sampling-and-voting and debate;
\cite{li2024,du2023}). Second, \textit{verification instead of trust}. Our manager
currently takes each worker's solved report and the notes at face value (in the v2
scaffold, it performs only a sanity check using existing sample cases provided by the
problem). A confidently-wrong DP, or a degenerate notes entry like the \texttt{abc399\_e}
loop (\S4.2), propagates unchecked and can overwrite a correct intermediate answer. Adding
a verifier that runs each candidate against a comprehensive suite of generated
tests (or, for mathematics, checks the argument with proof-verification software) before the
manager accepts it would catch precisely the regressions we observe, at the cost of extra
calls.

\subsection{Limitations}

\begin{itemize}
  \item \textbf{Serving-provider reliability (OpenRouter), \S2.4 only.} The five models of
  \S2.4 were run through OpenRouter's multi-provider routing, which proved a
  substantial and \textit{time-varying} confound. This is the limitation that motivated the
  pinned-backend set of \S2.1: those seven arms each talk to exactly one backend and are free
  of everything in this item, which is also why they - not the broader \S2.4 sweep - carry
  the paper's headline claims. Individual providers intermittently
  stalled (no response within multi-minute wall-clock caps), dropped the response
  mid-stream (\texttt{IncompleteRead}), returned error objects instead of a completion
  (missing choices, 5xx/504s), clamped output below the requested cap (e.g.\
  truncating generations at 32k while advertising 262k), or returned
  reasoning-only replies with \texttt{content=null}. When no provider yielded a
  usable completion, the attempt was recorded with no code and scored as a failure.
  Reported pass@1 therefore counts infrastructure failures as wrong answers, making
  absolute levels a conservative lower bound. Exact numbers are not perfectly
  reproducible, since they depend partly on provider health at run time. The relative
  single-vs-manager comparison is less sensitive to this confound because both conditions 
  use the same provider pool, although provider-level variation cannot be ruled out.

  \item \textbf{Provider-served weights may vary (\S2.4 only).} OpenRouter does not
  guarantee a single quantization or build per model; a config pinned to one provider set
  can be served differently than another, adding a second, smaller source of run-to-run
  variation on top of decoding temperature. The \S2.1 arms are pinned to one backend each,
  and the single locally served arm (Qwen3.8-27B) to a specific FP8 checkpoint, so this
  does not apply to them.

  \item \textbf{Single pass in \S2.4.} The thinking-on/off comparisons in \S2.4 are one
  pass per condition (to limit API cost), so per-model significance there rests on the 16k
  five-pass runs; those deltas should be read as point estimates within the
  ${\sim}4.5$\,pp pass-to-pass variation the 16k five-pass runs show on a
  manager-minus-single delta (individual arms vary by ${\sim}2.7$\,pp). This is the other
  reason \S2.1 repeats every condition five times, and it is why Opus-5's $85 \to 91$ -
  the highest score in the paper - remains a single-pass point estimate rather than a
  measured effect.

  \item \textbf{Fable 5 has no manager arm.} It was run single-only, so the strongest
  single-call result in the paper is also the one condition where we cannot say what the
  scaffold would do.

  \item \textbf{Refusals are scored as failures, and only one arm can incur them.} Fable 5
  returned \texttt{stop\_reason=refusal} on 6 of 500 problem-passes, spread sporadically
  over three ordinary problems (\S2.3). We deliberately configure no fallback model, since
  a cell labelled Fable 5 must contain Fable 5's outcome, but the consequence is that its
  $87.4$ carries ${\sim}1.2$ points of classifier penalty that no other arm in the paper is
  exposed to, and that the refusals add variance rather than a constant offset. Comparisons
  \textit{against} Fable 5 are therefore mildly conservative in its disfavour.

  \item \textbf{Qwen3.8's 128k single arm is a replay, not an independent run.} Those
  generations were produced at 250k and truncated to 128k after the fact (\S3.2). It is the
  right comparison for isolating the cap - identical generations either side - but it cannot
  capture any way the model might have budgeted its reasoning differently had it been told
  the smaller limit up front. A model that knows it has 128k may stop exploring sooner and
  write its answer earlier, in which case the replay understates what a real 128k run would
  score; we report the as-generated 250k reading alongside it throughout so the difference
  is visible rather than assumed.

  \item \textbf{A 9B model with thinking on was untestable.} Qwen3.5-9B with reasoning
  enabled could not be evaluated in any configuration we tried. Through OpenRouter it
  generates unbounded reasoning and returns reasoning-only replies
  (\texttt{content=null}, \texttt{finish=error}) or truncates before an answer appears,
  and no token or effort setting fixed it; every attempt was archived rather than graded.
  Locally the blocker was different: the checkpoint refused batching on ollama, leaving
  one sequence at a time and zero completed calls in a 34-minute smoke test. We report it
  as a model limitation rather than a data point, and it bounds how small a ``thinking''
  model this scaffold can be applied to.

  \item \textbf{One benchmark family.} All reported results are competitive-programming
  code generation. Math and knowledge benchmarks (AIME, MATH-500, GPQA, HLE) were run in
  exploratory form but are not reported here: on a probe of the ten hardest problems in
  each, the frontier models sit at or near ceiling (Opus-5 100\% on AIME, GPQA and
  MATH-500), leaving little headroom in which to measure a scaffold effect.
\end{itemize}

\section{Conclusion}

Holding the underlying model fixed, a lightweight manager--worker scaffold over a shared 
workspace with no training or task-specific tuning improves all three models measured with
both arms on a pinned backend and the verifier-gated v2 scaffold, over five paired passes
each: Qwen3.8-27B by $+23.4$, GPT-5.6-Luna by $+10.6$ and GPT-5.6-Terra by $+8.0$ points. The Qwen3.8 result is
the most striking: a 27B open-weight model
with the scaffold reaches $86.4$, level with the best single-call result in the study
(Claude Fable 5, $87.4 \pm 1.1$, run without any scaffold) and above Opus-5's single-call
$85$. It also lifts Opus-5 from $85$ to $91$ with thinking on, the highest absolute score
observed in the study, although this single-pass result warrants replication. With reasoning off,
the scaffold improves three of the four open models by $+3$ to $+42$ points, with the 
largest gains for Minimax-M3 and Kimi-K3. At the tight 16k cap, these models also show 
substantial reductions in unusable output, while Kimi-K3's $+42$-point gain at 128k occurs 
with zero recorded truncations in either arm. The results are conditional, however: 
Qwen3.6-35B shows no significant improvement at 16k and a $-9$-point change at 128k with
reasoning off. These results
provide a baseline against which more complex systems involving heterogeneous models or code
verification can be compared.

What the scaffold costs is the other half of the result. A manager roughly triples the bill, but in many cases it is still cheaper than switching to a larger model. With the manager, GPT-5.6-Terra lands 2.4 points short of Fable 5, unresolved at $p = 0.59$, for \$11.71 a pass against \$61.11. Qwen3.8-27B with a manager lands 1.0 point short, with a much lower parameter count suitable for local inference. The same trade shows up where GPT-5.6-Luna behind a manager matches GPT-5.6-Terra's unaided single call at 44\% of the
price, having started 9.8 points behind it. On this evidence the scaffold may be a cost effective way of improving results.

% IEEE reference style: numbered [n], initials before surname, article title in sentence
% case inside quotes, venue italicised, year last. thebibliography emits its own
% \section*{References}, so there is no heading above it. Entries are ordered by first
% citation in the text, which is what makes the numbers run 1, 2, 3 down the page: LaTeX
% numbers a \bibitem by its position in this list, not by where it is cited. Reordering
% entries therefore renumbers the paper, and \cite keeps every reference in step by
% itself. arXiv identifiers are kept on every e-print because arXiv asks for them in the
% reference list, and they must stay on one line with no markup inside the identifier.
\addcontentsline{toc}{section}{References}
\begin{thebibliography}{99}

\bibitem{guo2024} T. Guo, X. Chen, Y. Wang, R. Chang, S. Pei, N.~V. Chawla, O. Wiest, and
X. Zhang, ``Large language model based multi-agents: A survey of progress and
challenges,'' arXiv:2402.01680, 2024.

\bibitem{wangq2024} Q. Wang, Z. Wang, Y. Su, H. Tong, and Y. Song, ``Rethinking the bounds
of LLM reasoning: Are multi-agent discussions the key?'' in \textit{Proc. ACL}, 2024,
arXiv:2402.18272.

\bibitem{tran2026} D. Tran and D. Kiela, ``Single-agent LLMs outperform multi-agent systems
on multi-hop reasoning under equal thinking token budgets,'' arXiv:2604.02460, 2026.

\bibitem{yao2022} S. Yao, J. Zhao, D. Yu, N. Du, I. Shafran, K. Narasimhan, and Y. Cao,
``ReAct: Synergizing reasoning and acting in language models,'' in \textit{Proc. ICLR},
2023, arXiv:2210.03629.

\bibitem{shinn2023} N. Shinn, F. Cassano, E. Berman, A. Gopinath, K. Narasimhan, and
S. Yao, ``Reflexion: Language agents with verbal reinforcement learning,'' in
\textit{Proc. NeurIPS}, 2023, arXiv:2303.11366.

\bibitem{madaan2023} A. Madaan, N. Tandon, P. Gupta, S. Hallinan, L. Gao, S. Wiegreffe,
U. Alon, N. Dziri, S. Prabhumoye, Y. Yang, S. Gupta, B.~P. Majumder, K. Hermann,
S. Welleck, A. Yazdanbakhsh, and P. Clark, ``Self-refine: Iterative refinement with
self-feedback,'' in \textit{Proc. NeurIPS}, 2023, arXiv:2303.17651.

\bibitem{sakana2026} Sakana AI, ``Sakana Fugu technical report,'' arXiv:2606.21228, 2026.

\bibitem{nielsen2025} S. Nielsen, E. Cetin, P. Schwendeman, Q. Sun, J. Xu, and Y. Tang,
``Learning to orchestrate agents in natural language with the Conductor,'' in
\textit{Proc. ICLR}, 2026, arXiv:2512.04388.

\bibitem{ong2024} I. Ong, A. Almahairi, V. Wu, W.-L. Chiang, T. Wu, J.~E. Gonzalez,
M.~W. Kadous, and I. Stoica, ``RouteLLM: Learning to route LLMs with preference data,'' in
\textit{Proc. ICLR}, 2025, arXiv:2406.18665.

\bibitem{hong2023} S. Hong, M. Zhuge, J. Chen, X. Zheng, Y. Cheng, C. Zhang, J. Wang,
Z. Wang, S.~K.~S. Yau, Z. Lin, L. Zhou, C. Ran, L. Xiao, C. Wu, and J. Schmidhuber,
``MetaGPT: Meta programming for a multi-agent collaborative framework,'' in
\textit{Proc. ICLR}, 2024, arXiv:2308.00352.

\bibitem{qian2023} C. Qian, W. Liu, H. Liu, N. Chen, Y. Dang, J. Li, C. Yang, W. Chen,
Y. Su, X. Cong, J. Xu, D. Li, Z. Liu, and M. Sun, ``ChatDev: Communicative agents for
software development,'' in \textit{Proc. ACL}, 2024, arXiv:2307.07924.

\bibitem{wu2023} Q. Wu, G. Bansal, J. Zhang, Y. Wu, B. Li, E. Zhu, L. Jiang, X. Zhang,
S. Zhang, J. Liu, A.~H. Awadallah, R.~W. White, D. Burger, and C. Wang, ``AutoGen:
Enabling next-gen LLM applications via multi-agent conversation,'' in \textit{Proc. COLM},
2024, arXiv:2308.08155.

\bibitem{li2023} G. Li, H.~A.~A.~K. Hammoud, H. Itani, D. Khizbullin, and B. Ghanem,
``CAMEL: Communicative agents for `mind' exploration of large language model society,'' in
\textit{Proc. NeurIPS}, 2023, arXiv:2303.17760.

\bibitem{huang2023} D. Huang, J.~M. Zhang, M. Luck, Q. Bu, Y. Qing, and H. Cui,
``AgentCoder: Multi-agent-based code generation with iterative testing and
optimisation,'' arXiv:2312.13010, 2023.

\bibitem{islam2024} M.~A. Islam, M.~E. Ali, and M.~R. Parvez, ``MapCoder: Multi-agent code
generation for competitive problem solving,'' in \textit{Proc. ACL}, 2024,
arXiv:2405.11403.

\bibitem{xu2025} J. Xu, Q. Sun, P. Schwendeman, S. Nielsen, E. Cetin, and Y. Tang,
``TRINITY: An evolved LLM coordinator,'' in \textit{Proc. ICLR}, 2026, arXiv:2512.04695.

\bibitem{wei2026} M. Wei, X. Chen, X. Niu, and S. Chen, ``ARIADNE: Agentic reward-informed
adaptive decision exploration via blackboard-driven MCTS for competitive program
generation,'' arXiv:2605.02431, 2026.

\bibitem{han2025} B. Han and S. Zhang, ``Exploring advanced LLM multi-agent systems based
on blackboard architecture,'' arXiv:2507.01701, 2025.

\bibitem{du2023} Y. Du, S. Li, A. Torralba, J.~B. Tenenbaum, and I. Mordatch, ``Improving
factuality and reasoning in language models through multiagent debate,'' in
\textit{Proc. ICML}, 2024, arXiv:2305.14325.

\bibitem{li2024} J. Li, Q. Zhang, Y. Yu, Q. Fu, and D. Ye, ``More agents is all you need,''
\textit{Trans. Mach. Learn. Res.}, 2024, arXiv:2402.05120.

\bibitem{wangj2024a} J. Wang, J. Wang, B. Athiwaratkun, C. Zhang, and J. Zou,
``Mixture-of-Agents enhances large language model capabilities,'' in \textit{Proc. ICLR},
2025, arXiv:2406.04692.

\bibitem{wangj2024b} J. Wang, S. Jain, D. Zhang, B. Ray, V. Kumar, and B. Athiwaratkun,
``Reasoning in token economies: Budget-aware evaluation of LLM reasoning strategies,'' in
\textit{Proc. EMNLP}, 2024, arXiv:2406.06461.

\bibitem{openai2026price} OpenAI. ``API pricing.'' Accessed: Aug. 25, 2026. [Online].
Available: \url{https://developers.openai.com/api/docs/pricing}. Short-context tier; the
long-context tier and the cached-input discount are not applied here.

\bibitem{anthropic2026price} Anthropic. ``Pricing.'' Accessed: Aug. 25, 2026. [Online].
Available: \url{https://platform.claude.com/docs/en/about-claude/pricing}. Base input and
output rates; the Batch API discount and the prompt-caching multipliers are not applied
here.

\bibitem{openrouter2026} OpenRouter. ``Qwen3.8-27B.'' Accessed: Aug. 25, 2026. [Online].
Available: \url{https://openrouter.ai/qwen/qwen3.8-27b}. The hosted list rate the local
vLLM run is priced at.

\bibitem{qwen2026c} Qwen Team. ``Qwen3.8-27B model card.'' Hugging Face. Accessed: Aug.
25, 2026. [Online]. Available: \url{https://huggingface.co/Qwen/Qwen3.8-27B}. 27B parameters;
context window 262{,}144 tokens, confirmed against the FP8 checkpoint as served here
(vLLM reports \texttt{max\_model\_len}~$=$~262144 for \texttt{Qwen/Qwen3.8-27B-FP8}).

\bibitem{anthropic2026effort} Anthropic. ``Effort,'' ``Thinking'' and ``Task budgets.''
Accessed: Aug. 25,
2026. [Online]. Available: \url{https://platform.claude.com/docs/en/build-with-claude/effort},
\url{https://platform.claude.com/docs/en/build-with-claude/thinking} and
\url{https://platform.claude.com/docs/en/build-with-claude/task-budgets}. \texttt{effort}
takes \texttt{low}, \texttt{medium}, \texttt{high}, \texttt{xhigh} or \texttt{max}, and
``the API default is \texttt{high}'' - setting it explicitly ``produces exactly the same
behavior as omitting the \texttt{effort} parameter entirely.'' On Fable 5 ``thinking is
already on: no configuration needed,'' \texttt{display} defaults to \texttt{omitted}, and
``no \texttt{display} setting returns the raw chain of thought.'' Reasoning tokens ``are
billed as output tokens, even when the thinking text isn't returned to you, and they count
toward \texttt{max\_tokens} alongside the response text'' - it is ``a hard limit on total
output, thinking plus response text.'' On the cap itself: \texttt{task\_budget} is what
makes a budget visible (``the model sees a running countdown''), and it is advisory,
whereas ``the enforced limit on total output tokens is still \texttt{max\_tokens}, which
truncates the response.''

\bibitem{jain2024} N. Jain, K. Han, A. Gu, W.-D. Li, F. Yan, T. Zhang, S. Wang,
A. Solar-Lezama, K. Sen, and I. Stoica, ``LiveCodeBench: Holistic and contamination free
evaluation of large language models for code,'' arXiv:2403.07974, 2024.

\bibitem{qwen2026a} Qwen Team. ``Qwen3.5: Towards native multimodal agents.'' Accessed:
Aug. 17, 2026. [Online]. Available: \url{https://qwen.ai/blog?id=qwen3.5} and
\url{https://huggingface.co/Qwen/Qwen3.5-9B}. Published February 2026; Qwen3.5-9B: 9B
parameters.

\bibitem{qwen2026b} Qwen Team. ``Qwen3.6-35B-A3B model card.'' Hugging Face. Accessed:
Aug. 17, 2026. [Online]. Available:
\url{https://huggingface.co/Qwen/Qwen3.6-35B-A3B}. Published April 2026; 35B total
parameters, 3B activated per token.

\bibitem{minimax2026} MiniMax. ``MiniMax-M3.'' Accessed: Aug. 19, 2026. [Online].
Available: \url{https://www.minimax.io/blog/minimax-m3} and
\url{https://huggingface.co/MiniMaxAI/MiniMax-M3}. Released 1 June 2026; ${\sim}428$B
total parameters, ${\sim}23$B activated per token.

\bibitem{moonshot2026} Moonshot AI. ``Kimi K3.'' Accessed: Aug. 19, 2026. [Online].
Available: \url{https://www.kimi.ai/blog/kimi-k3} and
\url{https://platform.kimi.ai/docs/guide/kimi-k3-quickstart}.
${\sim}2.8$T total parameters, Mixture-of-Experts with 896 experts and ${\sim}16$
activated per token.

\bibitem{anthropic2026} Anthropic. ``Claude models overview.'' Accessed: Aug. 14, 2026.
[Online]. Available:
\url{https://platform.claude.com/docs/en/about-claude/models/overview}

\bibitem{openai2026} OpenAI. ``GPT-5.6 model documentation.'' Accessed: Aug. 19, 2026.
[Online]. Available: \url{https://developers.openai.com/api/docs/guides/latest-model} and
\url{https://developers.openai.com/api/docs/models/gpt-5.6-luna}. Three variants
(Sol, Terra, Luna), closed weights, parameter counts not published.

\end{thebibliography}

\end{document}
